%=======================================================================
%  ABSTRACT
%=======================================================================
\chapter*{Abstract}
\addcontentsline{toc}{chapter}{Abstract}

Most anomaly detectors on multivariate sensor streams watch one quantity: the size of a deviation. Relational faults defeat that. Every variable stays inside its nominal range while the relationship between them breaks, so the detector sees nothing, and when it does fire it cannot say which mechanism failed.

Causal discovery infers from observational data alone which variables directly drive which others, and returns a graph whose edges stand for mechanisms and not for correlations. The literature proposes it as a remedy for both difficulties. The supporting evidence is narrow. A published evaluation usually exercises one discovery algorithm, and usually takes the causal relations to be linear; where several algorithms are compared, the dataset, the downstream detector and the thresholding rule commonly move alongside the graph, so a reported gain cannot be attributed to the causal structure rather than to everything that moved with it.

This thesis holds fixed everything that usually moves with the graph. A controlled grid of 27 configurations crosses three causal discovery algorithms, PCMCI, GES and CAM-UV, one from each of the constraint-based, score-based and functional families, with three regression variants that relax the linearity assumption in turn: linear, cubic polynomial, and a random Fourier approximation to a radial basis function kernel. The three domains differ deliberately in scale and in physical character: humanoid robot telemetry at 244 monitored variables, the PTB-XL electrocardiogram corpus at 12, and the Tennessee Eastman process at 52. Everything downstream of the discovered graph stays constant, and 17 automated checks confirm that on every build instead of taking it on trust. Three tuned deep autoencoders and a parameter-free mean absolute $z$-score control run through the identical protocol.

On detection the causal pipeline is competitive without being uniformly best, and the character of the fault decides which way a domain falls. Where faults are relational, on the electrocardiogram data, the best causal configuration reaches an AUC-ROC of 0.8114 against 0.7077 for the best tuned neural baseline and 0.6136 for the parameter-free control. Where faults are amplitude excursions, on the industrial process, the causal pipeline trails the neural baselines, 0.9151 against 0.9997, and cannot be separated from the parameter-free control at 0.9286, which falls inside its confidence interval. One pattern runs through both. The causal advantage falls monotonically as the control's score rises, which makes the condition under which structure pays measurable before either method is built.

The case for causal structure therefore rests on explanation, and this thesis measures that case on two numbers instead of arguing it. Relevance asks how often the variables a method ranks highest are the ones the anomaly actually involved. Ground truth for it comes from outside the pipeline: the perturbed subsystem on the robot, the standard lead territory of each diagnosis on the electrocardiogram, the documented actuator on the plant. Consistency asks whether independent configurations agree on those variables. Both sit beside the rate an uninformed method would reach, because the share of variables that count as relevant differs by more than an order of magnitude between the three domains, and the raw percentages rank them backwards.

Measured this way the attributions carry real information on two domains of three. On the plant the causal algorithms select the documented actuator at eight to thirteen times the rate of random selection, and on the robot the nine configurations agree with one another at twenty-six times chance. The electrocardiogram reads as a failure until the classes are separated. Pooled, the figure sits at chance. Grouped by diagnosis, infarctions localise above chance in all three algorithms, and an infarction is defined by the territory its changes appear in; conduction and rhythm classes do not. The method localises what is anatomically localisable and fails where the notion of a correct lead is itself weaker.

The comparison against post-hoc explanation does not favour the causal pipeline. SHAP over the neural baselines was scored on the same classes and the same ground truth, on all three datasets. Paired over anomaly classes, only the electrocardiogram supplies enough scorable classes to decide anything: there SHAP beats PCMCI at the depth this thesis states its attribution results at, and again at the next depth down, and is indistinguishable from CAM-UV. It is indistinguishable from GES too at every depth at which GES is still ranking rather than returning the whole of the smaller pool it searches. The robot's three attack classes and the plant's five documented actuators cannot reach the 5\% level whatever the difference, so no verdict is reported on either. On the plant the causal method names the documented actuator in 13 of 15 configurations against 12, an ordering that reverses at five variables and vanishes at ten. On the electrocardiogram CAM-UV and SHAP rank the same twelve leads and sit level, at 0.95 times chance each, and neither localises anything. Comparing only where both sides search a pool of the same size, intrinsic attribution was therefore nowhere more relevant than post-hoc attribution, which is the opposite of what motivated the work and is reported as measured. The test adds a qualification rather than a consolation: on the one domain with enough ground truth to test, intrinsic attribution is measurably behind only under PCMCI and indistinguishable from post-hoc attribution otherwise, while the other two domains decide nothing in either direction. Three attack classes, five documented actuators and fourteen localisable diagnoses are the whole of what three public benchmarks supply, and two of those three quantities are too small to support a test at all. One qualification belongs with it: the methods do not all search the same number of variables, and because lift divides by the size of that search space it rewards the larger one, so the counts appear beside it throughout.

On consistency the causal side wins once, on the electrocardiogram, and the index alone undersells the win. Separate agreement about a class from agreement a method would reach whatever the class, and the causal attributions carry a margin of 0.267 there against 0.047 for the neural ones, whose agreement is close to a fixed ranking of the leads. That is also the dataset where causal relevance is weakest. The configurations agree with one another, specifically and about the record in front of them, while agreeing on leads no likelier than chance to be diagnostic. Agreement is not correctness. The case for reading an explanation off the model instead of fitting a second one to approximate it rests on what the explanation is, and it does not rest on the explanation scoring higher.

Attribution names the variable whose modelled relationship broke, not the one that deviated most, identifying the tablet sensor downstream of an LED attack in 9 of 9 robotics configurations and a joint temperature in 8 of 9 for a mechanical fault. Six of those nine came from graphs holding no LED channel at all: a correct explanation produced from a structure that does not contain it.

Two findings qualify the approach, and neither was sought. Detection performance cannot validate a discovered graph, since a graph carrying no observed-parent edges detected competitively. And treating the subsampling rate as a property of each algorithm forecloses cross-algorithm significance testing, leaving 54 of 108 within-domain comparisons formally unpairable.

\vspace{1em}
\noindent\textbf{Keywords:} causal discovery, anomaly detection, explainable AI, PCMCI, GES, CAM-UV, recursive least squares, Tennessee Eastman, PTB-XL, cyber-physical systems.
